%auto-ignore
\documentclass[main.tex]{subfiles}

\begin{document}

\section{Subprogram and Subprogram Independence}

\label{sec:subprogramIndependence}

To prove \cref{thm:FreeIndependencePrinciple}, we need to introduce a notion of \emph{subprograms} and then the \emph{independence} of subprograms.
Our key lemma in this section is \cref{lem:uncorrelationIndependence}, which shows that, if a subprogram is sufficiently uncorrelated from the previous subprogram, then it is also independent from it.
\begin{defn}
A program can be written down formally as a sequence of lines, each assigning a new vector generated from \refNonlin{} or \refMatMul{} to a new variable, like the examples above. A \emph{subprogram} of a program is a span of such lines. Given a program $\pi$ and subprograms $\pi_{1},\ldots,\pi_{k}$, we write $\pi=\pi_{1}|\cdots|\pi_{k}$ if $\pi$ consists of $\pi_{1}$ followed by $\pi_{2}$ followed by $\pi_{3}$, and so on, ending with $\pi_{k}$. Each vector in $\pi$ is introduced in a unique subprogram $\pi_{i}$, for which we write $v\in\pi_{i}$. By convention, all initial vectors in $\mathcal V$ are included in the first subprogram $\pi_1$.

For example, the simple \netsort{} program represented by \cref{eqn:MLP} can be split into two subprograms, one which introduces $x^l, h^l$ for $l=1,\ldots, R$, and another with $l=R+1,\ldots, L$.
As another example, \cref{eqn:forwardGIABreak} and \cref{eqn:backwardGIABreak} are two subprograms that comprise the program expressing the forward and backward propagations of the network.
In general, forward propagation and backward propagation form natural subprograms, where the latter depends on (the preactivations computed in) the former.

Next, we introduce the notion of subprogram independence, which generalizes the main property of Simple GIA Check (\cref{assm:simpleGIACheck}).
Before that, we first recall the notation for distributional equality.
\paragraph{Notations}
Given two random variables $X, \Zz$, and a $\sigma$-algebra $\Aa$, the notation $X \disteq_\Aa \Zz$ means that for any integrable function $\phi$ and for any random varible $Z$ measurable on $\Aa$, $\EV \phi(X) Z = \EV \phi(\Zz)Z$.
We say that $X$ is distributed as (or is equal in distribution to) $\Zz$ conditional on $\Aa$.
In case $\Aa$ is the trivial $\sigma$-algebra, we just write $X \disteq \Zz$.
If $X$ and $Z$ agrees almost surely, then we write $X \aseq Z$.
\end{defn}

\begin{defn}
\label{def:subprogramIndep}Consider a program divided into two subprograms: $\pi=\pi_{1}|\pi_{2}$. For each $W\in\mathcal{W}$, we create an iid copy $\bar{W}$. Let $\bar{\pi}_{2}$ be the same as $\pi_{2}$ except that each matrix $W$ used in $\pi_{2}$ is replaced with $\bar{W}$ in $\bar{\pi}_{2}$. Let $x^{1},\ldots,x^{k}$ denote the vectors of $\pi_{1}$, $y^{1},\ldots,y^{p}$ denote the vectors of $\pi_{2}$, and $\bar{y}^{1},\ldots,\bar{y}^{p}$ denote the vectors of $\bar{\pi}_{2}$. We say \emph{$\pi_{2}$ is independent from $\pi_{1}$} if we have the distributional equality
\[
(Z^{x^{1}},\ldots,Z^{x^{k}},Z^{y^{1}},\ldots,Z^{y^{p}})\disteq(Z^{x^{1}},\ldots,Z^{x^{k}},Z^{\bar{y}^{1}},\ldots,Z^{\bar{y}^{p}})
\]
or equivalently, the conditional distributional equality
\[
(Z^{y^{1}},\ldots,Z^{y^{p}})\disteq_{Z^{x^{1}},\ldots,Z^{x^{k}}}(Z^{\bar{y}^{1}},\ldots,Z^{\bar{y}^{p}}).
\]
\end{defn}

For example, if $\pi_{1}$ expresses the forward computation of a network satisfying Simple GIA Check (\cref{assm:simpleGIACheck}), and $\pi_{2}$ expresses the backward computation, then we can assume $W^\trsp$ used in $\pi_2$ (backpropagation) is independent from $W$ used in $\pi_1$ (forward propagation).
This then implies $\pi_{2}$ is independent from $\pi_{1}$.
On the other hand, the subprogram given by \cref{eqn:backwardGIABreak} is \emph{not} independent from the subprogram given by \cref{eqn:forwardGIABreak}, as we calculated.

Note \cref{def:subprogramIndep} is purely about the structure of the subprograms (which is the only thing the random variables $Z^\bullet$ depend on), not the values of the vectors or matrices in the program.

We will talk about polynomials of matrices below, e.g.\ $WW^\trsp W^2 (W^\trsp)^2$ if $W$ is a square matrix.
\begin{defn}[Noncommutative Polynomial]
    Given a matrix $W \in \R^{m\times n}$, a \emph{(noncommutative) monomial} in $\{W, W^\trsp\}$ is a product $a X_1 X_2 \cdots X_k$ for some $a\in \R$, $k \ge 0$, where each $X_i \in \{W, W^\trsp\}$ (with appropriate matching shapes).
    Its degree is $k$.
    A \emph{(noncommutative) polynomial} in  $\{W, W^\trsp\}$ is a linear combination of such monomials.
    Its degree is the maximum degree of all of its nonzero monomials.
\end{defn}
\begin{rem}\label{rem:calculatingPolynomial}
    If $W$ and $v$ are a matrix and a vector of a program, and $A$ is a polynomial in $\{W,W^\trsp\}$, then one can naturally express the product $Av$ in \netsort{}.
    For example, if $A = W (W^\trsp)^2 W$, then $Av = u$ where
    \[
        v^1 = Wv,\quad
        v^2 = W^\trsp v^1,\quad
        v^3 = W^\trsp v^2,\quad
        u = W v^3.\]
    Below, when we say a program \emph{calculates $Av$}, we mean a program of this form.
\end{rem}

The following is the key workhorse of this section. It says that, in a program $\pi_{1}|\pi_{2}$, if $\pi_{2}$ 1) consists of only \refMatMul{} and 2) only depends on $\pi_{1}$ through a single vector $v$ \emph{uncorrelated} with most vectors of $\pi_{1}$, then $\pi_{2}$ is independent from $\pi_{1}$. As we will comment shortly, this gives another proof of the Gradient Independence Assumption when Simple GIA Check (\cref{assm:simpleGIACheck}) is satisfied.

\begin{lemma}[Independence from Uncorrelation]
\label{lem:uncorrelationIndependence}Consider a \netsort{} program $\pi=\pi_{1}|\pi_{2}$, and $\pi_{2}$ calculates $Av$ (in the sense of \cref{rem:calculatingPolynomial}) where $v$ is a vector of $\pi_{1}$ and $A$ is a polynomial in $\{W,W^{\trsp}\}$ for some $W\in\mathcal{W}$. If $\EV Z^{v}Z^{x}=\EV Z^{v}\hat{Z}^{y}=0$ for every $x,y\in\pi_{1}$ with $y=Wx$ or $y=W^{\trsp}x$, then $\pi_{2}$ is independent from $\pi_1$.
\end{lemma}

Note if $Wv$ or $W^{\trsp}v$ appears in $\pi_{1}$, then the premise implies $Z^{v}=0$ and the theorem becomes vacuously true. So to get nontrivial implications, $Wv$ and $W^{\trsp}v$ should not appear in $\pi_{1}$.
Given \cref{lem:uncorrelationIndependence}, we can obtain a more explicit form of $Z^{Av}$ by some straightforward calculation.
\begin{lemma}
\label{lem:uncorrelatedDecomp}In the setting of \cref{lem:uncorrelationIndependence}, we have
\[
Z^{Av}=\tau Z^{v}+S
\]
where $S$ is zero-mean and independent from $\{Z^{u}:u\in\pi_{1}\}$, and $\tau=\lim_{n\to\infty}\frac{1}{n}\tr A$.
\end{lemma}

\begin{proof}
By \cref{lem:uncorrelationIndependence}, we can assume that $W$ (where $A$ is a polynomial in $\{W,W^{\trsp}\}$) is independent from everything in $\pi_{1}$. Then a simple induction on the degree of $A$ in $W$ shows that 
\begin{equation}
Z^{Av}=\tau Z^{v}+S\label{eq:dontknowtau}
\end{equation}
 where $S$ is zero-mean and independent from $\{Z^{u}:u\in\pi_{1}\}$.

However, at this point, we don't know what $\tau$ would be. To understand $\tau$, we proceed as follows. Let $\bar{v}$ be sampled iid according to $\bar{v}_{\alpha}\sim Z^{v}$ (which can be constructed with a \netsort{} program\footnote{The vector $v$ can always be expressed as $\phi(g^{1},\ldots,g^{k})$ for some G-vars $g^{1},\ldots,g^{k}$ and some function $\phi:\R^{k}\to\R$. Then $\bar{v}$ can be constructed as $\phi(\bar{g}^{1},\ldots,\bar{g}^{k})$ where $\{\bar{g}_{\alpha}^{1},\ldots,\bar{g}_{\alpha}^{k}\}\sim\Gaus(\mu,\Sigma)$ with $\mu_{i}=\EV Z^{g^{i}}$ and $\Sigma_{ij}=\EV Z^{g^{i}}Z^{g^{k}}$.}), so that $Z^{\bar{v}}\disteq Z^{v}$. Then we can easily see from the definition of $Z$ that $Z^{A\bar{v}}=\tau Z^{\bar{v}}+\bar{S}$ where $\tau$ here is the same as $\tau$ in \cref{eq:dontknowtau} and $(Z^{\bar{v}},\bar{S})\disteq(Z^{v},S)$. Then by the Master Theorem,
\[
\tau\EV(Z^{\bar{v}})^{2}=\EV Z^{A\bar{v}}Z^{\bar{v}}=\lim_{n\to\infty}\frac{1}{n}\bar{v}^{\trsp}A\bar{v}=\lim_{n\to\infty}\frac{1}{n}\tr A\EV(Z^{\bar{v}})^{2}
\]
we get $\tau=\lim_{n\to\infty}\frac{1}{n}\tr A$ as desired.
\end{proof}

\paragraph{Extension to \netsortplus{} and Variable Dimensions}
Both \cref{lem:uncorrelatedDecomp,lem:uncorrelationIndependence} hold as stated for \netsortplus{} programs (see \cref{sec:netsortplus}) and programs with variable dimensions (see \cref{appendix:VarDim}).
Note that, in programs with variable dimensions, the $\tau$ in \cref{lem:uncorrelatedDecomp} vanishes if $v$ and $Av$ are in different CDCs.

\paragraph{Another Look at How Simple GIA Check Implies GIA}

The lemmas above give us another proof that the gradient independence assumption leads to correct calculations given Simple GIA Check (\cref{thm:GIA}).
To demonstrate \cref{lem:uncorrelationIndependence,lem:uncorrelatedDecomp} in action, before proving FIP, we give a proof of \cref{thm:GIA} for MLP.
\begin{proof}[Proof sketch for an MLP]
Consider the MLP $f(\xi)=W^{L+1}x^{L}(\xi)$ with input $\xi\in\R^{d}$ and output dimension $1$, where we recursively define, for $l=1,\ldots,L$, 
\begin{align*}
g^{l}(\xi) & =W^{l}x^{l-1}(\xi) & dx^{l-1}(\xi) & =W^{l\trsp}dg^{l}(\xi) \\
x^{l}(\xi) & =\phi(g^{l}(\xi)) & dg^{l}(\xi) & =\phi'(g^{l}(\xi))\odot dx^{l}(\xi).
\end{align*}
Here, the dimensionalities are $W^1 \in \R^{n \times d}; W^2, \ldots, W^L \in \R^{n \times n}; W^{L+1} \in \R^{1 \times n}$, and for all $l\in[L]$, $g^l, x^l, dx^l, dg^l \in \R^n$.

Let $\pi_{F}$ be the program that computes $g^l, x^l, l=1,\ldots, L$, on an input $\xi$.
Below, we abbreviate $g^l = g^l(\xi)$ and so on.
In this program, $\mathcal V = \{g^1, dx^L = \sqrt n W^{L+1}\}$,
\footnote{
    Here $dx^l$ should be interpreted as $\sqrt n \partial f/\partial x^l$, and likewise for $dg^l$.
}%
and $\mathcal W = \{W^2, \ldots, W^L\}$.%
\footnote{
    For concreteness (which will not matter much below): we sample $W^l_{\alpha\beta} \sim \Gaus(0, 1/n)$ for $l = 2, \ldots, L+1$, and suppose $g^1_\alpha \disteq \Gaus(0, 1)$, induced by appropriate sampling of $W^1$.
}

Now let's see how to apply \cref{lem:uncorrelatedDecomp,lem:uncorrelationIndependence} to show \cref{thm:GIA}.
The first step of the backward pass is $dg^{L}=\phi'(g^{L})\odot dx^{L}$. Let $\pi_{1}$ denote $\pi_{F}$ plus this line. Let $\pi_{2}$ be the next line $dx^{L-1}=W^{L\trsp}dg^{L}$. Then $\pi_{1},\pi_{2}$ along with $v=dg^{L}$ satisfies the condition of \cref{lem:uncorrelationIndependence}.
Indeed, $Z^{dg^L}$ is uncorrelated from $Z^{g^L}$ and $Z^{x^{L-1}} = \hat Z^{x^{L-1}}$ because $Z^{dg^L}$ is linear in $Z^{dx^L}$, which is by definition sampled independently from all $W^1, \ldots, W^L$.
Thus, by \cref{lem:uncorrelationIndependence}, we may \emph{rigorously} treat $W^{L\trsp}$ as independent from $W^{L}$ for the purpose of computing any quantity of the form of the \cref{thm:NetsorTMasterTheorem} for the program $\pi_{1}|\pi_{2}$.
In particular, by \cref{lem:uncorrelatedDecomp}, $Z^{dx^{L-1}} = \tau Z^{dg^L} + S$, where $S$ is zero-mean and independent from $\{Z^{u}:u\in\pi_{1}\}$, and $\tau=\lim_{n\to\infty}\frac{1}{n}\tr W^{L\trsp} = 0$, i.e.\ $Z^{dx^{L-1}} = S$.

Then we can repeat this reasoning inductively: $Z^{dg^{L-1}} = \phi'(Z^{g^{L-1}}) \cdot Z^{dx^{L-1}} = \phi'(Z^{g^{L-1}}) \cdot S$ is uncorrelated from $Z^{g^{L-1}}$ and $Z^{x^{L-2}} = \hat Z^{x^{L-2}}$ because $Z^{dg^{L-1}}$ is linear in $S$.
This lets us apply \cref{lem:uncorrelationIndependence,lem:uncorrelatedDecomp} to $\pi_2$ being the line $dx^{L-2}=W^{L-1\trsp}dg^{L-1}$, $\pi_1$ being all previous lines, and $v$ being $dg^{L-1}$.
So on and so forth.
\end{proof}

\subsection{Proving \texorpdfstring{\cref{lem:uncorrelationIndependence}}{Key Lemma}}

Here we prove \cref{lem:uncorrelationIndependence}.
It will use the following trivial but useful facts.
\begin{prop}
    \label{prop:hatZRandomSource}For any $x\in\pi$, $Z^{x}$ is measurable against the $\sigma$-algebra generated by $\{\hat{Z}^{y}:\text{G-var }y\in\pi\}$.
    .
\end{prop}
    
\begin{lemma}
    \label{lemma:uncorImpliesPartialMeanZero}Fix a matrix $W$. If $\EV Z^{v}\hat{Z}^{Wh}=0$ for all G-var $Wh$ introduced before $v$, then $\EV\partial Z^{v}/\partial\hat{Z}^{Wh}=0$ for all such $Wh$ as well.
    
\end{lemma}
\cref{prop:hatZRandomSource} follows from a simple inductive argument, and \cref{lemma:uncorImpliesPartialMeanZero} follows from \cref{eq:ExpectationPartialDer}.

\begin{proof}[Proof of \cref{lem:uncorrelationIndependence}]
    It suffices to prove the case when we assume $A=\prod_{i=1}^{p}W^{t_{i}},t_{i}\in\{1,\trsp\}$, can be expressed as a monomial of degree $p$ in $\{W,W^{\trsp}\}$. Let $\bar{W}$ be an iid copy of $W$ as in \cref{def:subprogramIndep} and likewise let $\bar{\pi}_{2}$, $y^{1},\ldots,y^{p}$, $\bar{y}^{1},\ldots,\bar{y}^{p}$ be as in \cref{def:subprogramIndep}. By our assumption, $\pi_{2}$ computes $y^{1}=W^{t_{1}}v$ and $y^{i}=W^{t_{i}}y^{i-1}$ for $i=2,\ldots,p$, and likewise $\bar{\pi}_{2}$ computes $\bar{y}^{1}=\bar{W}^{t_{1}}v$ and $\bar{y}^{i}=\bar{W}^{t_{i}}\bar{y}^{i-1}$ for $i=2,\ldots,p$.
    
    Define $\mathcal{X}$ to be the $\sigma$-algebra generated by $\{Z^{x}:x\in\pi_{1}\}$. Note that by \cref{prop:hatZRandomSource}, $\mathcal{X}$ is also the $\sigma$-algebra generated by $\{\hat{Z}^{x}:\text{G-var }x\in\pi_{1}\}$. Then we need to show
    \[
    (Z^{y^{1}},\ldots,Z^{y^{p}})\disteq_{\mathcal{X}}(Z^{\bar{y}^{1}},\ldots,Z^{\bar{y}^{p}}).
    \]
    We proceed by induction on $p$. The base case of $p=0$ is vacuously true.
    
    Suppose the induction hypothesis holds for $p=q-1$ 
    \[
    (Z^{y^{1}},\ldots,Z^{y^{q-1}})\disteq_{\mathcal{X}}(Z^{\bar{y}^{1}},\ldots,Z^{\bar{y}^{q-1}}).
    \]
    and we shall prove it for $p=q$. A simple inductive argument (as in our proof of the semicircle law; see \cref{sec:semicircleProofGeneralK}) shows that for each $i\in[q-1]$, we have 
    \begin{equation}
    Z^{\bar{y}^{i}}=\tau_{i}Z^{v}+\bar{S}_{i}\label{eq:barYDecomp}
    \end{equation}
    for some $\tau_{i}\in\R$ and some zero-mean $\bar{S_{i}}$ independent from $\mathcal{X}$ (and thus also independent from $Z^{v}$). %
    
    Since we have assumed $A$ to be a monomial in $\{W,W^{\trsp}\}$, we can suppose WLOG that $t_{q}=1$ so that $y^{q}=Wy^{q-1}$. By definition, $Z^{y^{q}}=\hat{Z}^{y^{q}}+\Zdot^{y^{q}}$ and $Z^{\bar{y}^{q}}=\hat{Z}^{\bar{y}^{q}}+\Zdot^{\bar{y}^{q}}$. We shall show
    \[
    (Z^{y^{1}},\ldots,Z^{y^{q-1}},\hat{Z}^{y^{q}},\Zdot^{y^{q}})\disteq_{\mathcal{X}}(Z^{\bar{y}^{1}},\ldots,Z^{\bar{y}^{q-1}},\hat{Z}^{\bar{y}^{q}},\Zdot^{\bar{y}^{q}})
    \]
    which would imply the IH for $p=q$.
    
    \textbf{Case $\hat{Z}^{y^{q}}$: }By definition, $\hat{Z}^{y^{q}}$ (resp. $\hat{Z}^{\bar{y}^{q}}$) is zero-mean, jointly Gaussian with $\{\hat{Z}^{Wx}:x,Wx\in\pi\}$ (resp. $\{\hat{Z}^{\bar{W}x}:x,\bar{W}x\in\bar{\pi}_{2}\}$), and independent from $\{\hat{Z}^{Qx}:x,Qx\in\pi\}$ for any $Q\ne W$ (resp. for any $Q\ne\bar{W}$). By \cref{prop:hatZRandomSource}, these facts fully determine the distributions of $Z^{y^{1}},\ldots,Z^{y^{q-1}},\hat{Z}^{y^{q}}$ and of $Z^{\bar{y}^{1}},\ldots,Z^{\bar{y}^{q-1}},\hat{Z}^{\bar{y}^{q}}$ conditioned on $\mathcal{X}$. \emph{We thus need to show that A) $\Cov(\hat{Z}^{y^{q}},\hat{Z}^{y^{i}})=\Cov(\hat{Z}^{\bar{y}^{q}},\hat{Z}^{\bar{y}^{i}})$ for all $i=1,\ldots,q-1$, and B) $\Cov(\hat{Z}^{y^{q}},\hat{Z}^{g})=\Cov(\hat{Z}^{\bar{y}^{q}},\hat{Z}^{g})$ for all G-var $g\in\pi_{1}$.}
    
    \textbf{A) }For any $i=2,\ldots,q-1$, we either have $y^{i}=Wy^{i-1}$ or $y^{i}=W^{\trsp}y^{i-1}$. 1) In the latter case, $\hat{Z}^{y^{i}}$ is independent from $\hat{Z}^{y^{q}}$, and likewise $\hat{Z}^{\bar{y}^{i}}$ is independent from $\hat{Z}^{\bar{y}^{q}}$, so trivially \emph{$\Cov(\hat{Z}^{y^{q}},\hat{Z}^{y^{i}})=\Cov(\hat{Z}^{\bar{y}^{q}},\hat{Z}^{\bar{y}^{i}})=0$. }2) In the former case,
    \begin{align*}
    \Cov(\hat{Z}^{y^{q}},\hat{Z}^{y^{i}}) & =\Cov(\hat{Z}^{Wy^{q-1}},\hat{Z}^{Wy^{i-1}})\\
     & =\sigma_{W}^{2}\EV Z^{y^{q-1}}Z^{y^{i-1}}=\sigma_{W}^{2}\EV Z^{\bar{y}^{q-1}}Z^{\bar{y}^{i-1}}\quad\text{by IH}\\
     & =\Cov(\hat{Z}^{\bar{y}^{q}},\hat{Z}^{\bar{y}^{i}})
    \end{align*}
    as desired. For $i=1$, we either have $y^{1}=Wv$ or $y^{1}=W^{\trsp}v$, and a similar logic shows $\Cov(\hat{Z}^{y^{q}},\hat{Z}^{y^{1}})=\Cov(\hat{Z}^{\bar{y}^{q}},\hat{Z}^{\bar{y}^{1}})$ as well.
    
    \textbf{B) }If $g\in\pi_{1}$ is a G-var with $g=Qx$ for some $Q\ne W$ and $x\in\pi_{1}$, then by definition, $\hat{Z}^{g}$ is independent from both $\hat{Z}^{y^{q}}$ and $\hat{Z}^{\bar{y}^{q}}$. Now suppose instead $g=Wx$ for some $x\in\pi_{1}$. Then
    \begin{align*}
    \Cov(\hat{Z}^{y^{q}},\hat{Z}^{g}) & =\sigma_{W}^{2}\EV Z^{y^{q-1}}Z^{x}=\sigma_{W}^{2}\EV Z^{\bar{y}^{q-1}}Z^{x}\quad\text{by IH}\\
     & =\sigma_{W}^{2}(\tau_{q-1}\EV Z^{v}Z^{x}+\EV\bar{S}_{q-1}Z^{x})=0\quad\text{by \cref{eq:barYDecomp}}\\
     & =\Cov(\hat{Z}^{\bar{W}\bar{y}^{q-1}},\hat{Z}^{Wx})=\Cov(\hat{Z}^{\bar{y}^{q}},\hat{Z}^{g}).
    \end{align*}
    Here $\EV\bar{S}_{q-1}Z^{x}=0$ because $\bar{S}_{q-1}$ is zero-mean and independent from $Z^{x}$ as $x\in\pi_{1}$, and $\EV Z^{v}Z^{x}=0$ by the premise of this theorem. 
    
    \textbf{Case $\Zdot^{y^{q}}$: }By definition, $\Zdot^{y^{q}}$ is a linear combination $\Zdot^{y^{q}}=\sigma_{W}^{2}\sum_{(y,x)\in\mathcal{P}}Z^{x}\EV\frac{\partial Z^{y^{q-1}}}{\partial\hat{Z}^{y}}$, where $\mathcal{P}$ is the set of all $y=W^{\trsp}x$ introduced before $y^{q}$ in $\pi$. Likewise, $\Zdot^{\bar{y}^{q}}$ is a linear combination $\Zdot^{\bar{y}^{q}}=\sigma_{W}^{2}\sum_{i\in\mathcal{I}}Z^{\bar{y}^{i-1}}\EV\frac{\partial Z^{\bar{y}^{q-1}}}{\partial\hat{Z}^{\bar{y}^{i}}}$ where $\mathcal{I}$ is the set of all $i\in[q-1]$ such that $\bar{y}^{i}=\bar{W}^{\trsp}\bar{y}^{i-1}$, and for convenience here we have set $y^{0}=\bar{y}^{0}=v$. Clearly, by IH, 
    \[
    (Z^{\bar{y}^{1}},\ldots,Z^{\bar{y}^{q-1}},\hat{Z}^{\bar{y}^{q}},\Zdot^{\bar{y}^{q}})\disteq_{\mathcal{X}}\left(Z^{y^{1}},\ldots,Z^{y^{q-1}},\hat{Z}^{y^{q}},\sigma_{W}^{2}\sum_{i\in\mathcal{I}}Z^{y^{i-1}}\EV\frac{\partial Z^{y^{q-1}}}{\partial\hat{Z}^{y^{i}}}\right).
    \]
    Note that $\{(y^{i},y^{i-1}):i\in\mathcal{I}\}\subseteq\mathcal{P}$. So it suffices to show that $\EV\frac{\partial Z^{y^{q-1}}}{\partial\hat{Z}^{y}}=0$ for all $(y,x)\in\mathcal{P}\setminus\{(y^{i},y^{i-1}):i\in\mathcal{I}\}$. Fix one such $y$, which must have been introduced in $\pi_{1}$.
    
    Now by \cref{eq:barYDecomp}, $\EV\frac{\partial Z^{y^{q-1}}}{\partial\hat{Z}^{y}}=\EV\frac{\partial Z^{\bar{y}^{q-1}}}{\partial\hat{Z}^{y}}=\EV\tau_{q-1}\frac{\partial Z^{v}}{\partial\hat{Z}^{y}}+\frac{\partial\bar{S}_{q-1}}{\partial\hat{Z}^{y}}$. Since $\bar{S}_{q-1}$ is zero-mean and independent from $\mathcal{X}\ni\hat{Z}^{y}$, we have $\EV\frac{\partial\bar{S}_{q-1}}{\partial\hat{Z}^{y}}=0$. At the same time, by \cref{lemma:uncorImpliesPartialMeanZero}, $\EV\frac{\partial Z^{v}}{\partial\hat{Z}^{y}}=0$ because $\EV Z^{v}\hat{Z}^{W^{\trsp}h}=0$ for all G-var $W^{\trsp}h\in\pi_{1}$ by the premise of this theorem. Combining these results together, we get $\EV\frac{\partial Z^{y^{q-1}}}{\partial\hat{Z}^{y}}=0$ as desired. %
\end{proof}

    
\section{Proving \texorpdfstring{Free Independence Principle}{Free Independence Principle}}
\label{sec:FIP}

\begin{proof}[Proof of \cref{thm:FreeIndependencePrinciple}]

For each $i=1,\ldots,t$, let $A^{i}$ be a polynomial in $\{W,W^{\trsp}\}$ for some $W\in\mathcal{W}$ or in $\{\Diag(x):x\in\pi\}$, such that consecutive $A^{i}$ and $A^{i+1}$ are always polynomials in different collections. (Here superscripts always denote indices).
Then we need to show
\[
\frac{1}{n}\tr(A^{t}-I\frac{1}{n}\tr A^{t})\cdots(A^{1}-I\frac{1}{n}\tr A^{1})\asto0.
\]
We apply the trace trick to re-express this as
\begin{equation}
\frac{1}{n}\EV v^{\trsp}(A^{t}-I\frac{1}{n}\EV u^{t\trsp}A^{t}u^{t})\cdots(A^{1}-I\frac{1}{n}\EV u^{1\trsp}A^{1}u^{1})v \asto 0
\label{eq:MixedMomentConverges}
\end{equation}
where expectation is taken over $v,u^{1},\ldots,u^{t}\sim\Gaus(0,I)$.
We will express this as a \netsortplus{} program and apply \cref{thm:NetsorT+MeanConvergenceLinearlyBounded} to show this convergence.
This \netsortplus{} program takes the form of $\pi|\pi'$ where $\pi$ is the program in the theorem statement defining $W\in\mathcal{W}$ and $x\in\pi$, and $\pi'$ computes $\frac{1}{n}v^{\trsp}(A^{t}-\frac{1}{n}u^{t\trsp}A^{t}u^{t})\cdots(A^{1}-\frac{1}{n}u^{1\trsp}A^{1}u^{1})v$, as we describe in more detail below%
\footnote{For concreteness (which won't matter in the proof below): The initial set of vectors $\mathcal{V}_{\pi|\pi'}$ is then the initial vectors of $\pi$ along with $\{v,u^{1},\ldots,u^{t}\}$, $\mathcal{V}_{\pi|\pi'}=\mathcal{V}_{\pi}\cup\{v,u^{1},\ldots,u^{t}\}$.
Likewise, the initial matrices $\mathcal W_{\pi |\pi'}$ are the same as those of $\pi$, $\mathcal W_{\pi}$.}.
.

Inside $\pi'$ , we let $v^{0}=v$ and compute $v^{i}$ as follows in subprograms of $\pi'$:
\begin{align*}
v^{i} & =(A^{i}-\frac{1}{n}u^{i\trsp}A^{i}u^{i})v^{i-1},\quad\forall i\in[t].
\end{align*}
Then the quantity in \cref{eq:MixedMomentConverges} is just $\frac{1}{n}v^{\trsp}v^{t}$. %

\textbf{Example} If $A^{i}=(W^{\trsp})^{2}W$ for some $W\in\mathcal{W}$ then we can unpack $v^{i}=(A^{i}-\frac{1}{n}u^{i\trsp}A^{i}u^{i})v^{i-1}$ into the subprogram

\begin{align*}
g^{1} & =Wv^{i-1}&
g^{2} & =W^{\trsp}g^{2}&
g^{3} & =W^{\trsp}g^{3}=A^{i}v^{i-1}&
h^{1} & =Wu^{i}\\
h^{2} & =W^{\trsp}h^{1}&
h^{3} & =W^{\trsp}h^{2}&
c & =\frac{1}{n}u^{i\trsp}h^{3}&
v^{i} & =g^{3}-c\cdot v^{i-1}.
\end{align*}
Here $c$ is computed using \refMoment{} and $v^i$ is computed using \refNonlinPlus{}, and everything else uses \refMatMul{}.
Similarly, if $A^{i}=\Diag(x)$ for some $x\in\pi$ with bounded coordinates, then we can unpack $v^{i}=(A^{i}-\frac{1}{n}u^{i\trsp}A^{i}u^{i})v^{i-1}$ into
\begin{align*}
g & =x\odot v^{i-1}&
h & =x\odot u^{i}&
c & =\frac{1}{n}u^{i\trsp}h&
v^{i} & =g-c\cdot v.
\end{align*}
Here $c$ is computed using \refMoment{} and everything else is computed using \refNonlinPlus{}.
Note that, in both examples, the nonlinearities involved satisfy the nonlinearity conditions of \cref{thm:NetsorT+MeanConvergenceLinearlyBounded}.
One can easily see this is true in general for the entire subprogram $\pi'$.
By induction on $t$, we prove the following claim, which would imply \cref{eq:MixedMomentConverges}.
\begin{claim}
\emph{For each $v^{i},i\ge1$, the associated random variable $Z^{v^{i}}$ is zero-mean and 1) independent from $\mathcal{X}^{i-1}\defeq\{Z^{x}:\text{\ensuremath{x} introduced before or is \ensuremath{v^{i-1}}\}}$ if $A^{i}$ is not diagonal or 2) uncorrelated from $\mathcal{\mathcal{X}}^{i-1}$ and $\hat{\mathcal{X}}^{i-1}\defeq\{\hat{Z}^{x}:\text{G-var \ensuremath{x} introduced before or is \ensuremath{v^{i-1}}\}}$ if $A^{i}$ is diagonal.}
\end{claim}

Here, ``uncorrelated'' means $\EV Z^{v^{i}}R=0$ for all $R\in\mathcal{X}^{i-1}\cup\hat{\mathcal{X}}^{i-1}$. Applying this claim to $v^{t}$, we get
\[
\EV_{v, u^1, \ldots, u^t} v^{\trsp}(A^{t}-u^{t\trsp}A^{t}u^{t})\cdots(A^{1}-u^{1\trsp}A^{1}u^{1})v=v^{\trsp}v^{t}\asto\EV Z^{v}Z^{v^{t}}=0
\]
by \cref{thm:NetsorT+MeanConvergenceLinearlyBounded}%
\footnote{
    While \cref{thm:NetsorT+MeanConvergenceLinearlyBounded} is stated for a program (and not a subprogram), its proof can be readily adapted to the subprogram case.
}
, as desired, assuming rank stability (\cref{assm:asRankStab}). We shall check rank stability after proving this claim.
\begin{proof}
We proceed by induction on $t$, starting with $t=1$. Our key tool is \cref{lem:uncorrelatedDecomp}.
\begin{description}
\item [{BaseCase:}] By the cyclic property of trace, we can WLOG suppose $A^{1}$ is a polynomial in $\{W,W^{\trsp}\}$ for some $W\in\mathcal{W}$. Furthermore, we can assume WLOG that $A^{1}$ is a monomial in $W$ and $W^{\trsp}$. Then $\frac{1}{n}u^{1\trsp}A^{1}u^{1}\asto\tau^{1}\defeq\lim_{n\to\infty}\frac{1}{n}\tr A^{1}$. by something? By \cref{lem:uncorrelatedDecomp}, $Z^{v^{1}}=Z^{A^{1}v}=\tau^{1}Z^{v}+S$ where $S$ is zero-mean and independent from $\mathcal{X}^{0}$. Then $Z^{v^{1}}=Z^{A^{1}v-\left(\frac{1}{n}u^{1\trsp}A^{1}u^{1}\right)v}=S$ satisfies the required zero-mean and independence property.
\item [{Induction:}]~
\begin{description}
\item [{NondiagonalCase}] By IH, $Z^{v^{i-1}}$ is uncorrelated from $\mathcal{X}^{i-2}\cup\hat{\mathcal{X}}^{i-2}$. So we can apply \cref{lem:uncorrelatedDecomp} in the same fashion as in the base case to obtain the desired result.
\item [{DiagonalCase:}] Suppose $A^{i}=\Diag(x)$ for some $x\in\pi$. By the non-consecutive assumption, $A^{i-1}$ is not diagonal. So IH tells us $Z^{v^{i-1}}$ is independent from $\mathcal{X}^{i-2}$. One can see easily that $Z^{v^{i}}=Z^{v^{i-1}}(Z^{x}-\EV Z^{x}).$ We need to prove $\EV Z^{v^{i}}R=0$ for all $R\in\mathcal{X}^{i-1}\cup\mathcal{\hat{X}}^{i-1}$. We divide into cases:
\begin{enumerate}
\item Suppose $R\in\mathcal{X}^{i-2}\cup\hat{\mathcal{X}}^{i-2}$. Then $\EV Z^{v^{i}}R=\EV Z^{v^{i-1}}(Z^{x}-\EV Z^{x})R=\EV Z^{v^{i-1}}\EV(Z^{x}-\EV Z^{x})R=0$ because $Z^{v^{i-1}}$ is independent from $Z^{x}$ and $R$ and is zero-mean.
\item Suppose $R\in\mathcal{\hat{X}}^{i-1}\setminus\mathcal{\hat{X}}^{i-2}$. Then $R$ is possibly correlated with $Z^{v^{i-1}}$ but $\{R,Z^{v^{i-1}}\}$ is independent from $\mathcal{\mathcal{X}}^{i-2}\ni x$ by \cref{lem:uncorrelationIndependence}. Therefore, $\EV Z^{v^{i}}R=\EV Z^{v^{i-1}}(Z^{x}-\EV Z^{x})R=\EV Z^{v^{i-1}}R\EV(Z^{x}-\EV Z^{x})=0$.
\item Suppose $R\in\mathcal{X}^{i-1}\setminus\mathcal{X}^{i-2}$. Then $R$ decomposes into $R=\tau Z^{v^{i-2}}+S$ by \cref{lem:uncorrelatedDecomp} for some $\tau\in\R$ and $S$ zero-mean and independent from $\mathcal{X}^{i-2}$. Then 
\begin{align*}
\EV Z^{v^{i}}R & =\EV Z^{v^{i-1}}=\EV Z^{v^{i-1}}(Z^{x}-\EV Z^{x})R\\
 & =\EV Z^{v^{i-1}}(Z^{x}-\EV Z^{x})(\tau Z^{v^{i-2}}+S)\\
 & =\tau\EV Z^{v^{i-1}}(Z^{x}-\EV Z^{x})Z^{v^{i-2}}+\EV Z^{v^{i-1}}(Z^{x}-\EV Z^{x})S\\
 & =0.
\end{align*}
Here $\EV Z^{v^{i-1}}(Z^{x}-\EV Z^{x})S=\EV Z^{v^{i-1}}S\EV(Z^{x}-\EV Z^{x})=0$ because $\{Z^{v^{i-1}},S\}$ is independent from $Z^{x}$. Similarly, $\EV Z^{v^{i-1}}(Z^{x}-\EV Z^{x})Z^{v^{i-2}}=\EV Z^{v^{i-1}}\EV(Z^{x}-\EV Z^{x})Z^{v^{i-2}}=0$ because by IH $Z^{v^{i-1}}$ is zero-mean and independent from $\mathcal{X}^{i-2}\ni Z^{x},Z^{v^{i-2}}$.
\end{enumerate}
\end{description}
\end{description}
\end{proof}

Finally, to see rank stability (\cref{assm:asRankStab}), we simply note two points: 1) The subprogram $\pi$ is a \netsort{} program and thus it has rank stability automatically (\cref{remk:rankStabilityNetsor}). 2) $\hat Z^x$ has nonzero variance for every G-var $x$ computed in $\pi'$.
Since each vector used in \refMatMulPlus{} in $\pi'$ depends linearly on some unique $\hat Z^x$, this shows that the rank in \cref{assm:asRankStab} is always full, and thus rank stability holds.

\end{proof}

    
\end{document}